\chapter{svnmucc Reference---Subversion Multiple URL Command
Client}\label{svn.ref.svnmucc}

The Subversion Multiple URL Command Client (\texttt{svnmucc}) is a tool
that can make arbitrary changes to the repository without the use of a
working copy. As regards remote commit capabilities, the functionality
provided by this tool is similar to, but far exceeds, that which is
offered by the Subversion command-line client itself. For example,
\texttt{svnmucc} is not limited to performing only a single type of
change in a given commit. It is also able to modify file content and
versioned properties without a working copy, which is functionality not
currently offered by \texttt{svn}.

This reference describes the \texttt{svnmucc} tool, and the various
remote modification actions you can perform using it.

\section{svnmucc}\label{svn.ref.svnmucc.re}

\emph{Perform one or more Subversion repository URL-based ACTIONs,
committing the result as a (single) new revision.}

\textbf{Synopsis}

\texttt{svnmucc\ ACTION...}

\subsection{Description}\label{svn.ref.svnmucc.re.desc}

\texttt{svnmucc} is a program for modifying Subversion-versioned data
without the use of a working copy. It allows operations to be performed
directly against the repository URLs of the files and directories that
the user wishes to change. Each invocation of \texttt{svnmucc} attempts
one or more \textless ACTION\textgreater s, atomically committing the
results of those combined \textless ACTION\textgreater s as a single new
revision.

\subsection{Actions}\label{svn.ref.svnmucc.re.actions}

\texttt{svnmucc} supports the following actions (and related arguments),
which may be combined into ordered sequences on the command line:

\begin{description}
\item[cp \textless REV\textgreater{} \textless SRC-URL\textgreater{}
\textless DST-URL\textgreater{}]
Copy the file or directory located at \textless SRC-URL\textgreater{} in
revision \textless REV\textgreater{} to \textless DST-URL\textgreater.
\item[mkdir \textless URL\textgreater{}]
Create a new directory at \textless URL\textgreater. The parent
directory of \textless URL\textgreater{} must already exist (or have
been created by a prior \texttt{svnmucc} action), as this command does
not offer the ability to automatically create any missing intermediate
parent directories.
\item[mv \textless SRC-URL\textgreater{}
\textless DST-URL\textgreater{}]
Move the file or directory located at \textless SRC-URL\textgreater{} to
\textless DST-URL\textgreater.
\item[rm \textless URL\textgreater{}]
Delete the file or directory located at \textless URL\textgreater.
\item[put \textless SRC-FILE\textgreater{} \textless URL\textgreater{}]
Add a new file---or modify an existing one---located at
\textless URL\textgreater, copying the contents of the local file
\textless SRC-FILE\textgreater{} as the new contents of the created or
modified file. As a special consideration,
\textless SRC-FILE\textgreater{} may be \texttt{-} to instruct
\texttt{svnmucc} to read from standard input rather than a local
filesystem file.
\item[propset \textless NAME\textgreater{} \textless VALUE\textgreater{}
\textless URL\textgreater{}]
Set the value of the property \textless NAME\textgreater{} on the target
\textless URL\textgreater{} to \textless VALUE\textgreater.
\item[propsetf \textless NAME\textgreater{} \textless FILE\textgreater{}
\textless URL\textgreater{}]
Set the value of the property \textless NAME\textgreater{} on the target
\textless URL\textgreater{} to the contents of the file
\textless FILE\textgreater.
\item[propdel \textless NAME\textgreater{} \textless URL\textgreater{}]
Delete the property \textless NAME\textgreater{} from the target
\textless URL\textgreater.
\end{description}

\subsection{Options}\label{svn.ref.svnmucc.re.sw}

Options specified on the \texttt{svnmucc} command line are global to all
actions performed by that command line. The following is a list of the
options supported by this tool:

\begin{description}
\item[\texttt{-\/-config-dir} \textless DIR\textgreater{}]
Read configuration information from the specified directory instead of
the default location (\texttt{.subversion} in the user\textquotesingle s
home directory).
\item[\texttt{-\/-config-option} \textless CONFSPEC\textgreater{}]
Set, for the duration of the command, the value of a runtime
configuration option. \textless CONFSPEC\textgreater{} is a string which
specifies the configuration option namespace, name and value that
you\textquotesingle d like to assign, formatted as
\textless FILE\textgreater:\textless SECTION\textgreater:\textless OPTION\textgreater={[}\textless VALUE\textgreater{]}.
In this syntax, \textless FILE\textgreater{} and
\textless SECTION\textgreater{} are the runtime configuration file
(either \texttt{config} or \texttt{servers}) and the section thereof,
respectively, which contain the option whose value you wish to change.
\textless OPTION\textgreater{} is, of course, the option itself, and
\textless VALUE\textgreater{} the value (if any) you wish to assign to
the option. For example, to temporarily disable the use of the
compression in the HTTP protocol, use
\texttt{-\/-config-option=servers:global:http-compression=no}. You can
use this option multiple times to change multiple option values
simultaneously.
\item[\texttt{-\/-extra-args} (\texttt{-X})
\textless ARGFILE\textgreater{}]
Read additional would-be command-line arguments from
\textless ARGFILE\textgreater, one argument per line. As a special
consideration, \textless ARGFILE\textgreater{} may be \texttt{-} to
indicate that additional arguments should be read instead from standard
input.
\item[\texttt{-\/-file} (\texttt{-F}) \textless MSGFILE\textgreater{}]
Use the contents of the \textless MSGFILE\textgreater{} as the log
message for the commit.
\item[\texttt{-\/-help} (\texttt{-h}, \texttt{-?})]
Show program usage information and exit.
\item[\texttt{-\/-message} (\texttt{-m}) \textless MSG\textgreater{}]
Use \textless MSG\textgreater{} as the log message for the commit.
\item[\texttt{-\/-no-auth-cache}]
Prevent caching of authentication information (e.g., username and
password) in the Subversion runtime configuration directories.
\item[\texttt{-\/-non-interactive}]
Disable all interactive prompting (e.g., requests for authentication
credentials).
\item[\texttt{-\/-revision} (\texttt{-r}) \textless REV\textgreater{}]
Use revision \textless REV\textgreater{} as the baseline revision for
all changes made via the \texttt{svnmucc} actions. This is an important
option which users should habituate to using whenever modifying existing
versioned items to avoid inadvertently undoing contemporary changes made
by fellow team members.
\item[\texttt{-\/-root-url} (\texttt{-U})
\textless ROOT-URL\textgreater{}]
Use \textless ROOT-URL\textgreater{} as a base URL to which all other
URL targets are relative. This URL need not be the
repository\textquotesingle s root URL (such as might be reported by
\texttt{svn\ info}). It can be any URL common to the various targets
which are specified in the \texttt{svnmucc} actions.
\item[\texttt{-\/-password} (\texttt{-p})
\textless PASSWD\textgreater{}]
Use \textless PASSWD\textgreater{} as the password when authenticating
against a Subversion server. If not provided, or if incorrect,
Subversion will prompt you for this information as needed.
\item[\texttt{-\/-username} \textless NAME\textgreater{}]
Use \textless USERNAME\textgreater{} as the username when authenticating
against a Subversion server. If not provided, or if incorrect,
Subversion will prompt you for this information as needed.
\item[\texttt{-\/-version}]
Display the program\textquotesingle s version information and exit.
\item[\texttt{-\/-with-revprop}
\textless NAME\textgreater=\textless VALUE\textgreater{}]
Set the value of the revision property \textless NAME\textgreater{} to
\textless VALUE\textgreater{} on the committed revision.
\end{description}

\subsection{Examples}\label{svn.ref.svnmucc.re.examples}

To (safely) modify a file\textquotesingle s contents without using a
working copy, use \texttt{svn\ cat} to fetch the current contents of the
file, and \texttt{svnmucc\ put} to commit the edited contents thereof.

\begin{verbatim}
$ # Calculate some convenience variables.
$ export FILEURL=http://svn.example.com/projects/sandbox/README
$ export BASEREV=`svn info ${FILEURL} | \
                  grep '^Last Changed Rev' | cut -d ' ' -f 2`
$ # Get a copy of the file's current contents.
$ svn cat ${FILEURL}@${BASEREV} > /tmp/README.tmpfile
$ # Edit the (copied) file.
$ vi /tmp/README.tmpfile
$ # Commit the new content for our file.
$ svnmucc -r ${BASEREV} put README.tmpfile ${FILEURL} \
          -m "Tweak the README file."
r24 committed by harry at 2013-01-21T16:21:23.100133Z
# Cleanup after ourselves.
$ rm /tmp/README.tmpfile
\end{verbatim}

Apply a similar approach to change a file or directory property. Simply
use \texttt{svn\ propget} and \texttt{svnmucc\ propsetf} instead of
\texttt{svn\ cat} and \texttt{svnmucc\ put}, respectively.

\begin{verbatim}
$ # Calculate some convenience variables.
$ export PROJURL=http://svn.example.com/projects/sandbox
$ export BASEREV=`svn info ${PROJURL} | \
                  grep '^Last Changed Rev' | cut -d ' ' -f 2`
$ # Get a copy of the directory's "license" property value.
$ svn -r ${BASEREV} propget license ${PROJURL} > /tmp/prop.tmpfile
$ # Tweak the property.
$ vi /tmp/prop.tmpfile
$ # Commit the new property value.
$ svnmucc -r ${BASEREV} propsetf prop.tmpfile ${PROJURL} \
          -m "Tweak the project directory 'license' property."
r25 committed by harry at 2013-01-21T16:24:11.375936Z
# Cleanup after ourselves.
$ rm /tmp/prop.tmpfile
\end{verbatim}

Let\textquotesingle s look now at some multi-operation examples.

To implement a ``moving tag'', where a single tag name is recycled to
point to different snapshots (for example, the current latest stable
version) of a codebase, use \texttt{svnmucc\ rm} and
\texttt{svnmucc\ cp}:

\begin{verbatim}
$ svnmucc -U http://svn.example.com/projects/doohickey \
          rm tags/latest-stable \
          cp HEAD trunk tags/latest-stable \
          -m "Slide the 'latest-stable' tag forward."
r134 committed by harry at 2013-01-12T11:02:16.142536Z
$ 
\end{verbatim}

In the previous example, we slyly introduced the use of the
\texttt{-\/-root-url\ (-U)} option. Use this option to specify a base
URL to which all other operand URLs are treated as relative (and save
yourself some typing).

The following shows an example of using \texttt{svnmucc} to, in a single
revision, create a new tag of your project which includes a newly
created descriptive file and which lacks a directory which
shouldn\textquotesingle t be included in, say, a release tarball.

\begin{verbatim}
$ echo "This is the 1.2.0 release." | \
       svnmucc -U http://svn.example.com/projects/doohickey \
               -m "Tag the 1.2.0 release." \
               -- \
               cp HEAD trunk tags/1.2.0 \
               rm tags/1.2.0/developer-notes \
               put - tags/1.2.0/README.tag
r164 committed by cmpilato at 2013-01-22T05:26:15.563327Z
$ svn log -c 164 -v http://svn.example.com/projects/doohickey
------------------------------------------------------------------------
r164 | cmpilato | 2013-01-22 00:26:15 -0500 (Tue, 22 Jan 2013) | 1 line
Changed paths:
   A /tags/1.2.0 (from /trunk:163)
   A /tags/1.2.0/README.tag
   D /tags/1.2.0/developer-notes

Tag the 1.2.0 release.
$ 
\end{verbatim}

The previous example demonstrates not only how to do several different
things in a single \texttt{svnmucc} invocation, but also the use of
standard input as the source of new file contents. Note the presence of
\texttt{-\/-} to indicate that no more options follow on the command
line. This is required so that the bare \texttt{-} used in the
\texttt{svnmucc\ put} action won\textquotesingle t be flagged as a
malformed option indicator.

